% capitolo 8: integrazione della business intelligence
\chapter{Integrazione con la Business Intelligence}
\label{cap:business-intelligence}

\intro{Questo capitolo mette in relazione i due temi studiati durante il percorso di stage, il catalogo dati e la Business Intelligence, delineando il flusso ideale e sicuro delle informazioni.}\\

\section{Introduzione alla Business Intelligence}
La Business Intelligence (BI) è una disciplina che combina infrastrutture tecnologiche, strumenti di analisi dei dati, visualizzazioni avanzate e \emph{best practice} 
organizzative con l'obiettivo di supportare le aziende nell'adozione di decisioni basate su evidenze empiriche ed oggettive (\emph{data-driven decision making})\cite{site:business-intelligence}.

Dal punto di vista teorico e operativo, il flusso della Business Intelligence si articola in quattro passaggi sequenziali fondamentali:
\begin{enumerate}
    \item \textbf{Raccolta ed estrazione}: i dati grezzi provenienti da fonti eterogenee (quali database, fogli di calcolo Excel o sistemi \emph{on-premise}) vengono estratti e centralizzati in depositi unificati;
    \item \textbf{Analisi descrittiva}: le informazioni memorizzate vengono elaborate per comprendere l'andamento storico e corrente dell'azienda. In questa fase si opera una distinzione essenziale rispetto alla Business Analytics (BA): mentre la Business Intelligence si concentra sull'analisi descrittiva (rispondendo alle domande \emph{cosa} e \emph{come}), la Business Analytics dà priorità all'analisi predittiva e prescrittiva, sfruttando modelli statistici per prevedere eventi futuri (rispondendo alla domanda \emph{perché});
    \item \textbf{Visualizzazione dei dati}: i risultati dell'elaborazione analitica vengono convertiti in rappresentazioni grafiche intuitive (istogrammi, diagrammi a torta, \emph{scatter plot}), rendendo le metriche complesse accessibili e leggibili anche per gli utenti di business non tecnici;
    \item \textbf{Azione basata su informazioni in tempo reale}: la visualizzazione delle performance storiche e attuali, calata nel contesto aziendale, consente ai \emph{decision maker} di intervenire prontamente per ottimizzare i processi operativi.
\end{enumerate}

\begin{figure}[H]
    \centering
    \includegraphics[width=1\textwidth]{bi/bi_concept.png}
    \caption[Flusso logico della Business Intelligence]{Rappresentazione del flusso logico della Business Intelligence. \textit{Nota: La presente infografica è stata generata con l'ausilio di un sistema di Intelligenza Artificiale}}
    \label{fig:bi_concept}
\end{figure}

All'interno dell'infrastruttura di Wintech, la piattaforma di riferimento per l'analisi aziendale e la Business Intelligence è Microsoft Power BI\cite{site:powerbi-docs}. 
Questa soluzione fornisce servizi e strumenti in grado di connettere, trasformare e condividere i dati dell'organizzazione, traducendoli in informazioni dettagliate e visualizzazioni grafiche.

Utilizzato quotidianamente dall'azienda, Power BI soddisfa nativamente numerosi requisiti legati all'osservabilità e alla gestione del dato, offrendo:
\begin{itemize}
    \item \textbf{Centralizzazione delle fonti}: la capacità di connettere in un unico ambiente un'ampia varietà di dati strutturati (come database SQL Server, PostgreSQL) e non strutturati (come file JSON, CSV, fogli Excel), fungendo da singola fonte di verità per la reportistica aziendale;
    \item \textbf{Sincronizzazione programmata}: l'allineamento costante e automatico delle informazioni con le fonti originali, tramite Power BI Service e i \emph{Gateway} dati locali;
    \item \textbf{Organizzazione logica}: la strutturazione dei modelli di dati attraverso prospettive distinte (vista Report, vista Tabella, vista Modello), che permettono di esplorare le relazioni tra le entità importate;
    \item \textbf{Sicurezza e gestione degli accessi}: l'implementazione di standard di sicurezza, includendo un sistema RBAC a livello di \emph{workspace} e la sicurezza a livello di riga (RLS - \emph{Row-Level Security}).
\end{itemize}

\subsection{La piattaforma Microsoft Power BI}
Wintech adotta Microsoft Power BI come soluzione di riferimento per lo sviluppo dei propri sistemi di reporting interni e per la consulenza offerta ai clienti. 
Si tratta di una piattaforma avanzata di Business Intelligence che fornisce una \emph{suite} di servizi, applicazioni e connettori che trasformano molteplici 
fonti di dati in informazioni strategiche, dettagliate e visivamente interattive.

La piattaforma offre tre ambienti operativi principali, che supportano l'analista durante l'intero ciclo di modellazione e fruizione dei dati:
\begin{itemize}
    \item \textbf{Power BI Desktop}: Applicazione client di sviluppo utilizzata per connettersi alle sorgenti, effettuare la pulizia e la modellazione dei dati tramite il Power Query Editor, definire le misure in linguaggio DAX (\emph{Data Analysis Expressions}) e creare i report visivi.
    \item \textbf{Power BI Service}: Servizio in cloud basato su modello SaaS (\emph{Software as a Service}) che permette la pubblicazione, condivisione e distribuzione sicura dei report all'interno dell'organizzazione, nonché la creazione di dashboard riepilogative destinate al \emph{management}.
    \item \textbf{Power BI Mobile}: Applicazioni per dispositivi mobili ottimizzate per visualizzare i report in mobilità, sfruttando specifici layout reattivi.
\end{itemize}

\begin{figure}[H]
    \centering
    
\includegraphics[width=0.25\textwidth]{logo-powerbi.png}
    \caption{Logo ufficiale di Microsoft Power BI}
    \label{fig:power_bi_logo}
\end{figure}

\subsection{Contesto aziendale}
In Wintech S.p.A., la Business Intelligence ha un ruolo strategico. In qualità di consulente e \emph{System Integrator}, l'azienda aiuta i clienti a integrare 
sistemi ERP e a creare dashboard personalizzate, che permettono di monitorare facilmente le vendite, la logistica, le risorse umane e l'andamento finanziario.

Durante lo stage, lo studio di Power BI ha permesso di comprendere l'utilità pratica dello strumento, sviluppando report basati sia sul database di test \emph{AdventureWorks} 
sia sui movimenti contabili reali estratti in ambiente sicuro dai sistemi gestionali aziendali.

A titolo esemplificativo, la Figura~\ref{fig:report_powerbi} mostra l'interfaccia di un report sviluppato sui dati di test, evidenziando le funzionalità di 
Data Visualization e il pannello per la configurazione degli avvisi di monitoraggio.

\begin{figure}[H]
    \centering
    \includegraphics[width=1\textwidth]{bi/report_powerbi.png}
    \caption[Interfaccia di un report in Microsoft Power BI]{Esempio di dashboard analitica sviluppata in Power BI sui dati di \emph{AdventureWorks}. Sulla destra è visibile il pannello di configurazione degli avvisi, utile per il monitoraggio parziale dei dati}
    \label{fig:report_powerbi}
\end{figure}

Tuttavia, Power BI è uno strumento progettato specificamente per la Data Visualization e la reportistica. Proprio per questo motivo, elabora le informazioni 
esattamente come vengono estratte, senza verificarne preventivamente la qualità o la presenza di dati personali alla sorgente. 

Affinché la Business Intelligence sia davvero affidabile e conforme alle normative sulla privacy, è necessario creare un flusso informativo intermedio. 
Questo livello si posiziona tra le fonti grezze e i report finali, affidando l'ispezione, la pulizia e la classificazione dei dati a piattaforme specifiche di Data Discovery e Data Governance.


\section{Flusso sicuro dei dati}
% PER SPIEGARE QUAL è L'IDEA DIETRO LA CONNESSIONE DATA CATALOG - POWER BI + grafico del flusso:  dati grezzi → catalogazione → classificazione PII → data quality e 
% monitoraggio → mascheramento → governance e policy → strumenti BI che consumano solo dati puliti, sicuri e certificati
L'efficacia strategica della Business Intelligence non dipende solo dalle capacità analitiche o grafiche delle dashboard, ma è condizionata alla radice dall'accuratezza e 
dalla sicurezza del flusso informativo sottostante. Nelle architetture tradizionali, il collegamento diretto tra gli strumenti di Business Intelligence (come Microsoft 
Power BI) e le basi dati genera criticità significative sotto il profilo operativo, della Data Quality e della conformità al GDPR.

Connettere gli strumenti di Business Intelligence direttamente alle sorgenti espone l'organizzazione a tre vulnerabilità principali:
\begin{itemize}
    \item \textbf{Esposizione incontrollata di dati sensibili}: i report di BI analizzano tipicamente metriche aggregate. Tuttavia, una connessione diretta ai database di produzione consente agli analisti di accedere e importare record contenenti informazioni personali in chiaro (quali nomi, indirizzi, e-mail o codici fiscali), violando i principi di minimizzazione e limitazione delle finalità previsti dal GDPR;
    \item \textbf{Propagazione di anomalie di Data Quality}: se una sorgente dati presenta valori nulli, duplicati o anomalie strutturali, lo strumento di BI importa passivamente tali errori. Questo genera visualizzazioni distorte o errate, compromettendo l'affidabilità delle metriche esposte e quindi dei processi decisionali;
    \item \textbf{Assenza di contesto semantico e documentale}: gli analisti faticano a individuare e comprendere i dataset d'interesse (Data Discovery), dovendo decifrare strutture fisiche complesse e nomenclature tecniche prive di una descrizione formalizzata o di un'associazione chiara con i termini di business dell'organizzazione.
\end{itemize}

Per mitigare tali rischi e implementare nativamente i principi di \emph{Privacy by Design}, l'architettura informativa deve integrare un livello intermedio che funga 
da \emph{gatekeeper} logico. Un Data Catalog come OpenMetadata si interpone tra i sistemi gestionali e le piattaforme analitiche, orchestrando il ciclo di 
vita dei metadati e assicurando che la Business Intelligence consumi esclusivamente risorse certificate e sicure.

Il processo che trasforma i dati grezzi in informazioni utili per le decisioni aziendali si articola in sei fasi:
\begin{enumerate}
    \item \textbf{Estrazione dei metadati}: il modulo di Data Ingestion del catalogo interroga periodicamente le sorgenti in modalità \emph{Read-Only}, estraendo la sola topologia strutturale delle tabelle (nomi, colonne, tipi di dato) e un campione ristretto di record di esempio (\emph{Sample Data}), senza duplicare fisicamente i dati reali o appesantire l'infrastruttura transazionale;
    \item \textbf{Arricchimento semantico}: i metadati tecnici vengono associati ai termini formalizzati nel glossario aziendale, definendo responsabilità e \emph{ownership} per ciascun asset e facilitando la Data Discovery da parte degli utenti di business;
    \item \textbf{Individuazione e classificazione dei PII}: mediante l'esecuzione dell'\emph{AutoClassification Agent} o l'applicazione manuale dei tag, il sistema rileva le informazioni personali e vi applica i tag pertinenti (es. \texttt{PII.Sensitive}), mappando la distribuzione dei dati protetti;
    \item \textbf{Monitoraggio della qualità}: il motore di Data Quality esegue i test pianificati sui database sorgente. In caso di anomalie, il sistema segnala l'errore sulla dashboard di Data Observability e notifica i responsabili tramite alert automatici, impedendo la propagazione dell'errore a valle;
    \item \textbf{Mascheramento e anonimizzazione}: nel catalogo l'accesso ai dati sensibili del \emph{Sample Data} viene protetto tramite le policy RBAC, oscurando i valori (\texttt{******}) ai profili non autorizzati. A livello di database fisico, i dataset destinati ad ambienti di test e reporting vengono opportunamente anonimizzati tramite tecniche e modelli formali, mantenendo i soli vincoli relazionali necessari alle query di analisi;
    \item \textbf{Consumo sicuro}: lo sviluppatore o l'analista di BI si connette esclusivamente alle sorgenti o alle aree di \emph{staging} sicure contenenti i dati precedentemente catalogati, verificati e bonificati, garantendo una reportistica attendibile, accurata e conforme alle normative.
\end{enumerate}

\begin{figure}[H]
    \centering
    \includegraphics[width=1\textwidth]{bi/flusso_dati.png}
    \caption{Rappresentazione logica del flusso sicuro delle informazioni}
    \label{fig:flusso_dati}
\end{figure}

\subsection{Vincoli di integrazione}
Durante lo stage, a causa di vincoli legati alla disponibilità di licenze Microsoft Power BI Pro e alle policy di rete interne necessarie per l'esposizione delle API, 
non è stato possibile collaudare la connessione fisica tra l'istanza locale di OpenMetadata e Power BI. 
Tali requisiti, di competenza dell'area IT Infrastructure, sono stati posticipati per motivi di sicurezza perimetrale, portando lo stage a concentrarsi sullo studio 
concettuale e sulla modellazione logica del flusso informativo.

Un altro vincolo tecnico consiste nell'asimmetria dell'integrazione tra i due sistemi. All'interno di OpenMetadata è possibile registrare Power BI come un servizio sorgente 
per estrane i metadati, importando nel catalogo i metadati relativi a dashboard, report, datamodel e Data Lineage.

Tuttavia, questo flusso di metadati è esclusivamente unidirezionale:
\begin{itemize}
    \item Power BI alimenta il data catalog fornendo la topologia delle proprie risorse analitiche;
    \item Power BI non può in alcun modo connettersi a OpenMetadata per interrogarlo come fonte dati reale o per estrarre direttamente tabelle fisiche pulite e anonimizzate.
\end{itemize}

Questo limite deriva dall'architettura del catalogo dati: la piattaforma non opera come database fisico e non fa una copia delle informazioni aziendali, ma si limita 
ad archiviare e indicizzare i metadati descrittivi e qualitativi degli asset esterni. Di conseguenza, la pulizia e il mascheramento fisico dei dati devono essere delegati 
ad altri strumenti, che operano a livello di base dati o staging area, lasciando a OpenMetadata il solo ruolo di tracciamento e controllo di conformità.



\section{Strumenti di Data Discovery}
Durante l'analisi condotta nel percorso di stage, un quesito strategico fondamentale ha guidato la valutazione dell'architettura aziendale: è strettamente indispensabile 
introdurre una piattaforma di data catalog complessa e centralizzata come OpenMetadata, o è possibile soddisfare i requisiti di Data Governance combinando le funzionalità 
di Microsoft Power BI con strumenti verticali e leggeri di PII Discovery?

Power BI, ampiamente integrato nel perimetro operativo di Wintech S.p.A., si attesta come un eccellente strumento di reportistica e visualizzazione, ma implementa nativamente 
anche alcune funzionalità tipiche di un catalogo dati:
\begin{itemize}
    \item \textbf{Data Ingestion:} si connette a una vastissima gamma di fonti dati tramite connettori predefiniti;
    \item \textbf{Data Discovery:} permette l'esplorazione dei modelli semantici e la ricerca visiva dei dataset;
    \item \textbf{Data Collaboration:} supporta la condivisione sicura dei report attraverso l'integrazione con Microsoft Teams e Microsoft 365;
    \item \textbf{Data Observability:} consente il monitoraggio dei KPI e l'attivazione di allarmi in tempo reale al variare delle metriche aziendali.
\end{itemize}

Tuttavia, la business intelligence da sola non è in grado di garantire il controllo qualitativo e la conformità normativa alla sorgente. Power BI non effettua test di 
Data Quality automatici sui database sorgente e non dispone di algoritmi per scovare le informazioni personali nei database relazionali aziendali prima che queste vengano 
importate nel modello. 

Per colmare questo divario senza incorrere nella complessità sistemistica di un data catalog completo, lo stage ha previsto la valutazione di cinque soluzioni specialistiche 
(\emph{tool} di PII Discovery e anonimizzazione), analizzandone l'utilità, le caratteristiche e la fattibilità di integrazione nell'infrastruttura aziendale:
\begin{enumerate}
    \item \textbf{Microsoft Presidio}: un \emph{framework} \emph{open-source} e \emph{stateless} sviluppato da Microsoft, specializzato nell'individuazione e nell'anonimizzazione delle PII tramite modelli NLP locali (spaCy) ed espressioni regolari. Poiché non si connette direttamente alle fonti dati ma riceve testo tramite API, risulta estremamente flessibile per essere integrato all'interno di pipeline Python personalizzate prima del caricamento in Power BI;
    \item \textbf{ARX (Data Anonymization Tool)}: uno strumento \emph{open-source} dedicato alla de-identificazione fisica dei dati strutturati. Si distingue per la capacità di connettersi ai database relazionali, analizzare preventivamente il rischio statistico di re-identificazione e applicare modelli matematici formali di privacy, quali la \emph{k-Anonymity}, la \emph{l-Diversity} e la \emph{t-Closeness}, esportando tabelle bonificate e sicure utilizzabili dagli analisti di business intelligence;
    \item \textbf{Hawk-Eye}: una potente utility a riga di comando \emph{open-source} progettata per scansionare sia fonti strutturate (database SQL) sia risorse non strutturate (file PDF, Word, immagini e video) grazie all'uso combinato di pattern matching e tecnologie OCR. Pur eccellendo nella fase di Data Discovery e auditing, Hawk-Eye presenta due limiti significativi: non supporta nativamente Microsoft SQL Server e non applica alcuna tecnica di mascheramento, delegando la protezione a passaggi successivi dello sviluppo;
    \item \textbf{Bytebase}: uno strumento focalizzato sul DevSecOps e sulla gestione del ciclo di vita dei database. Consente l'applicazione di policy di mascheramento dinamico (\emph{dynamic data masking}) direttamente a livello di query, ma richiede l'adozione di un'infrastruttura di gestione complessa;
    \item \textbf{Fides (di Ethyca)}: una piattaforma di privacy engineering basata sulla metodologia \emph{Privacy-as-Code}, progettata per automatizzare il rispetto del GDPR direttamente nei flussi di sviluppo software. Richiede tuttavia un'installazione e una manutenzione sistemistica non trascurabile tramite container Docker.
\end{enumerate}

La combinazione di Power BI con un applicativo di PII Discovery leggero e mirato (come l'accoppiata ARX e Microsoft Presidio) rappresenta un'alternativa concreta ed 
efficiente per mappare e proteggere le informazioni sensibili dell'azienda. Questa architettura ibrida permette di blindare la sicurezza del dato alla radice e di 
abilitare un reporting sicuro, rinunciando tuttavia ad alcune funzionalità esclusive di una piattaforma di Data Governance nativa, quali la gestione di un glossario 
aziendale centralizzato per gli utenti di business, la tracciabilità visiva del Data Lineage e gli strumenti di collaborazione diretta (apertura di task, thread di discussione 
e gestione degli incidenti qualitativi).



\section{Considerazioni sull'adozione}
A completamento del \emph{Proof of Concept} sviluppato in ambiente locale, la scelta strategica sulla futura evoluzione dello stack tecnologico di Wintech S.p.A. 
deve fondarsi su un bilancio oggettivo tra i benefici operativi e l'onere di adozione, manutenzione e gestione infrastrutturale delle diverse soluzioni analizzate.

Il primo parametro discriminante è rappresentato dai requisiti hardware e dalla complessità di installazione e manutenzione delle piattaforme di catalogazione:
\begin{itemize}
    \item \textbf{OpenMetadata}: richiede l'esecuzione simultanea di quattro container Docker (il database transazionale MySQL/PostgreSQL, il server centrale in Java, il motore di ingestion basato su Apache Airflow e il cluster di indicizzazione Elasticsearch/OpenSearch), con un consumo minimo garantito di 6-8 GB di RAM e 4 vCPU per ambienti di produzione\cite{site:openmetadata-docs, site:manuale-utente};
    \item \textbf{Open Data Discovery (ODD)}: si attesta come un candidato \emph{open-source} alternativo più leggero, richiedendo solo tre componenti (la piattaforma ODD, un database PostgreSQL e l'agente collector), con requisiti sistemistici ridotti a 4 GB di RAM e 2 vCPU. Tuttavia, non dispone di un motore nativo di Auto-Classification per i dati sensibili, vanificando parte dei vantaggi della catalogazione automatica;
    \item \textbf{Microsoft Purview}: rappresenta il modello di riferimento enterprise per l'ecosistema Microsoft aziendale, integrando nativamente l'intelligenza artificiale e la Data Compliance in \emph{cloud}, ma comporta costi di licenza e complessità di configurazione (richiedendo spesso l'uso di Kubernetes) sproporzionati rispetto alle reali necessità correnti dell'organizzazione;
    \item \textbf{Tool leggeri standalone (ARX, Presidio, Hawk-Eye)}: non richiedono un'infrastruttura server sempre attiva (24/7), potendo essere eseguiti come script Python locali schedulati periodicamente in background (tramite Windows Task Scheduler o cron in Linux), riducendo a zero l'onere sistemistico fisso e azzerando i costi di adozione.
\end{itemize}

Sulla base di queste considerazioni, l'introduzione immediata di una piattaforma di Data Catalog completa come OpenMetadata risulterebbe parzialmente sovrodimensionata per 
l'attuale contesto operativo di Wintech S.p.A., determinando un rischio di sovra-ingegnerizzazione (\emph{over-engineering}) a fronte di un patrimonio informativo che può 
essere protetto con strategie più snelle.

Tuttavia, il vero valore del progetto di stage non risiede nell'adozione immediata e vincolante di un singolo applicativo software, bensì nella maturità e nella consapevolezza 
metodologica acquisite dall'azienda in materia di Data Governance, qualità del dato e Sviluppo Sicuro. Lo studio compiuto ha evidenziato come le tematiche di conformità al 
GDPR (inclusa la corretta classificazione dei tag e l'esigenza di una tracciabilità trasparente dei flussi) non siano meri vincoli burocratici, ma requisiti architetturali 
ad alto valore aggiunto che coinvolgono trasversalmente molteplici reparti aziendali:
\begin{itemize}
    \item Il team di Ricerca e Sviluppo acquisisce strumenti concreti per bonificare e anonimizzare i database di produzione prima di utilizzarli negli ambienti locali di test e sviluppo software;
    \item Il reparto Business Intelligence ottiene linee guida chiare per progettare flussi informativi intermediati e sicuri, riducendo il rischio di esporre accidentalmente record riservati o di propagare anomalie qualitative nelle metriche direzionali;
    \item L'area Compliance e Sicurezza dispone di un \emph{framework} di riferimento (ispirato agli standard ISO/IEC 27001 e ISO/IEC 29100) per gestire con efficacia gli audit interni, mappare i rischi di re-identificazione statistica e rispondere tempestivamente alle richieste di esercizio dei diritti degli interessati (come il diritto all'oblio).
\end{itemize}

In conclusione, lo scenario evolutivo più realistico e sostenibile per l'organizzazione prevede l'adozione di un approccio modulare e progressivo: anziché mantenere attiva 
un'infrastruttura di catalogazione complessa a livello aziendale, Wintech S.p.A. potrà capitalizzare i risultati dello stage implementando in produzione gli strumenti di 
PII Discovery e anonimizzazione deterministica (come ARX) per la bonifica dei database fisici operazionali, garantendo il pieno rispetto della privacy fin dalla progettazione 
(\emph{Privacy by Design}) all'interno di flussi analitici e applicativi intrinsecamente sicuri.